\section{Discussion Summary}

\begin{itemize}
    \item Focus primarily on refining the type denotation and semantics, postponing concrete syntax issues until later.
    \item Leverage the perspective of higher-order symbolic execution to formalize the approach, aiming to provide theoretical guarantees and clarify how our previous examples and programs map to this framework.
    \item Develop a more intuitive and user-friendly mechanism compared to traditional forward search strategies.
    \item Explicitly address support for recursion in the type system and reasoning principles.
\end{itemize}

\section{Intuition}

Design an underapproximate type system that is rich enough to solve programs from different domains—some interesting problems may require effects, higher-order, or relational features.

\section{Examples}

\paragraph{Trace Contracts} Numerous examples can be drawn from the literature on trace contracts and higher-order contracts. For instance, in dataflow analysis frameworks, proving algorithm termination often hinges on showing that the transfer function is monotonic. Beyond termination, one may also establish \emph{soundness} (i.e., that all output labels are reachable) and \emph{completeness} (all reachable labels can appear as output). 
Although many contract examples are implemented in Racket...
% they can be systematically translated into OCaml with the assistance of modern LLM tools.

\paragraph{Completeness of Algorithms} A question is whether we can formally establish the completeness of standard algorithms. As a case in point, breadth-first search (BFS) is considered complete, even over infinite graphs—provided there exists a finite reachable path, BFS will eventually discover it. Conversely, depth-first search (DFS) is not complete, as it may traverse an infinite branch and never reach some nodes. Dijkstra's algorithm, by contrast, is both sound and complete under standard assumptions.

\paragraph{Optimization and Reachability} Program optimization frequently relies on reasoning about reachability: is a particular code location necessarily (must) or possibly (may) reachable? Proving (un)reachability is central to dead code elimination and related analyses.

\paragraph{Incorrectness} Many existing examples in intermediate languages focus on unsafe behaviors, although most of them are about memory safety violations.

\paragraph{Unreachability and Divergence} Showing that a program may or must diverge (fail to terminate) is valuable for verifying liveness properties or for identifying infinite loops that might block further progress in a system.

\paragraph{Security} We can reuse the IFC example. However, many security guarantees are inherently relational. On the other hand, taint analysis can be another example to consider (as far as I know, it is not sound).

\section{Type}

\newcommand{\divergeTy}[1]{\Code{diverge}^{#1}}
\newcommand{\errTy}[1]{\Code{err}^{#1}}

\begin{figure}[t!]
    {\small
    \begin{alignat*}{2}
        \\\text{\textbf{Base Types}}& \quad & b  ::= \quad & \Int ~|~ \Bool ~|~ \Unit ~|~ ...
        \\\text{\textbf{Basic Types}}& \quad & s  ::= \quad & b ~|~ t \sarr t
        \\\text{\textbf{Refinement Types}}& \quad & t  ::= \quad & \divergeTy{s} ~|~ \errTy{s} ~|~ \rt{b}{\phi} ~|~
        \\ & \quad & \quad & \ol{x}{:}t\sarr t ~|~ \ul{x}{:}t\sarr t ~|~ \ol{x}{:}b\garr t ~|~ \ul{x}{:}b\garr t ~|~ t \unionty t ~|~ t \interty t
        \\\text{\textbf{Type Contexts}}& \quad &\Gamma ::= \quad  &\emptyset ~|~ \ol{x}{:}t, \Gamma ~|~ \ul{x}{:}t, \Gamma
      \end{alignat*}
    }
    \vspace*{-.15in}
    \caption{\DSL{} types}
    \label{fig:type-syntax}
    \vspace*{-.15in}
    \end{figure}

\paragraph{Type Syntax} Now the new types are more close to dependent types, however I believe there should still be some restrictions.

\paragraph{Examples}
\begin{enumerate}
    \item Base coverage type: \\
    $\urt{b}{\phi} \equiv r{:}b\olgarr \rt{b}{\phi \land \vnu = r}$.
    \item Function coverage type: \\
    $x{:}\ort{b_1}{\phi_1}\sarr\urt{b_2}{\phi_2} \equiv \ol{x}{:}\rt{b_1}{\phi_1}\sarr(\ol{r}{:}b_2\garr\rt{b_2}{\phi_2 \land \vnu = r})$.
    \item Higher-order coverage type: \\
    $y{:}(x{:}\ort{b_1}{\phi_1}\sarr\urt{b_2}{\phi_2})\sarr t \equiv \ol{y}{:}(\ol{x}{:}\rt{b_1}{\phi_1}\sarr(\ol{r}{:}b_2\garr\rt{b_2}{\phi_2 \land \vnu = r})) \sarr t$.
    \item Incorrectness triple: \\
    $[P(x)] y = f(x) [Q(y)] \equiv \ol{r}{:}b\garr \rt{b}{P(\vnu)} \ularr \rt{b}{Q(\vnu) \land \vnu = r}$. 
    \item (New!) bound constraints \\
    $\ul{r}{:}b\garr\rt{b}{\vnu > r}$: the term has a lower bound.
    \item (New!) Existential higher parameter, higher-order symbolic execution \\
    $\ul{f}{:}( ...) \sarr \Unit $: there exists a function $f$ that satisfies the given constraints can make the term terminate.
\end{enumerate} 

\paragraph{Type denotation} Assume there is no randomness in the language.
\begin{align*}
    \Code{Terminate}(e) &\doteq \exists v. \mstep{e'}{v} \lor \mstep{e}{\Code{err}} \\
    \denotation{ \errTy{s} } &\doteq \{ e ~|~ \emptyset \basicvdash e : s \land \mstep{e}{\Code{err}} \} \\
    \denotation{ \divergeTy{s} } &\doteq \{ e ~|~ \emptyset \basicvdash e : s \land \neg \Code{Terminate}(e) \} \\
    \denotation{ \rt{b}{\phi} } &\doteq \{ e ~|~ \emptyset \basicvdash e : b \land \Code{Terminate}(e) \land \forall v. \mstep{e}{v} \implies \phi[\vnu \mapsto v] \} \\
    \denotation{ \ol{x}{:}t_x\sarr t } &\doteq \{ e ~|~ \emptyset \basicvdash e : ... \land \forall v\in\denot{t_x}. e\ v \in \denot{t[x\mapsto v]} \} \\
    \denotation{ \ul{x}{:}t_x\sarr t } &\doteq \{ e ~|~ \emptyset \basicvdash e : ... \land \exists v\in\denot{t_x}. e\ v \in \denot{t} \} \\
    \denotation{ \ol{x}{:}b\garr t } &\doteq \{ e ~|~ \emptyset \basicvdash e : ... \land \forall v{:}b. e \in \denot{t[x\mapsto v]} \} \\
    \denotation{ \ul{x}{:}b\garr t } &\doteq \{ e ~|~ \emptyset \basicvdash e : ... \land \exists v{:}b. e \in \denot{t[x\mapsto v]} \} \\
    \denotation{ t \unionty t } &\doteq \denot{t} \cup \denot{t'} \\
    \denotation{ t \interty t } &\doteq \denot{t} \cap \denot{t'}
\end{align*}

I don't believe the underapproximate parameter should be dependent, in a reverse hoare logic / IL, we say for all outputs, there exists inputs -- if the constraints of output is dependent on inputs values, it hard to express underapproximation. And I don't find any benefit of dependent underapproximate parameter. There are many difficulties that allowing a under variable to be refered by both term and its type.


\paragraph{Typing} Variable.

\begin{align*}
    \ul{x}{:}\rt{b}{\phi}  &\vdash x : ? \\
    \ol{x}{:}\rt{b}{\phi}  &\vdash x \typeinfer y{:}b\olgarr \rt{b}{\phi \land \vnu = y} \quad \text{the inferred type is independent from the $x$}\\
    \ul{x}{:}\rt{b}{\phi}  &\vdash x \typeinfer \rt{b}{\vnu = x}\\
\end{align*}